\section{A spherical grid with an entropy balance}
\label{l:grid:section}

This argument records a cell containing the signal after each block of
samples.  The outgoing memory state can make the posterior concentrate
inside that cell.  We therefore choose a finer cell from the new posterior
and include its entire label in the analysis history.  The decrease in
cell size pays for that additional information.  A projection estimate
then bounds the remaining information acquired in one block.

The use of independent points and inverse distances to their successive
affine spans is related to the projection method of Sharan, Sidford, and
Valiant~\cite[Section~7, Lemma~12]{SSV2019}.  We prove the estimates used
here for exact Gaussian labels.  The choice of cell below is a global
maximum over all finer levels; it need not be a stopping time along the
grid.

We use the information conventions of Section~\ref{l:model}, and write
$h$ for differential entropy when it is finite.  Let $\sigma$ be uniform
probability on $S^{d-1}$, and take $d$ larger than an absolute constant.  Set
\begin{equation}\label{l:grid:parameters}
 k=p=\lfloor d/8\rfloor,\qquad q=d/2,\qquad \bar q=q+2.
\end{equation}
In particular, $q<d-1$ and $q-(p-1)-k$ is a positive multiple of $d$.
Constants denoted by $C$ are positive and absolute and may increase.

The cap estimate in Lemma~\ref{l:caps-endpoint} gives, for
$z\in S^{d-1}$ and $r>0$,
\begin{equation}\label{l:grid:cap}
 \sigma(B(z,r))\le r^{d-1},
\end{equation}
where $B(z,r)$ is the closed Euclidean ball.

For every integer $j\ge0$, let $\delta_j=2^{-j}$ and partition
$\mathbb R^d$ into half-open cubes of side $\delta_j/\sqrt d$, using
the same origin at every level.  These partitions are nested.  Write
$Q_j(s)$ for the unique cube containing $s$.  Its diameter is
$\delta_j$, so \eqref{l:grid:cap}, centered at any point of its
intersection with the sphere, gives
\begin{equation}\label{l:grid:cell-mass}
 \sigma(Q_j(s))\le 2^{-(d-1)j}.
\end{equation}
At most $(C2^j)^d$ level-$j$ cubes meet the sphere.  Indeed, their union
lies in the ball of radius $1+\delta_j\le2$, and their interiors are
disjoint.  The volume $v_d$ of the unit ball satisfies
$v_d\le(2\pi e/d)^{d/2}$: integrate $e^{-\|x\|^2/2}$ over the ball
of radius $\sqrt d$ and compare with its full Gaussian integral.
Dividing the volume of the radius-two ball by the cube volume now gives
the asserted count.

\Needspace{6\baselineskip}
\begin{lemma}[A refinement whose label is charged]
\label{l:grid:refinement}
Let $S\sim\sigma$, and let $G,W$ be countable discrete variables with
$H(G,W)<\infty$.  Suppose $G$ determines an integer $a\ge0$ and a
level-$a$ cube containing $S$ almost surely, and that $\mathbb E a<\infty$.
For every pair $(g,w)$ of positive probability, put
$\mu=\mathcal L(S\mid G=g,W=w)$.  For $\mu$-almost every $s$, define
\[
 J(s)=\max_{j\ge a}
 \bigl\{\log\mu(Q_j(s))+q(j-a)\log2\bigr\},
\]
and let $b(s)$ be the least level attaining this maximum.  In the joint
law define
\[
 N=\lfloor J(S)/d\rfloor,\qquad Z=(b(S),Q_{b(S)}(S)).
\]
The choices are measurable and finite almost surely.  They satisfy
$\mathbb E b<\infty$, $H(G,W,Z)<\infty$, and
\begin{align}
 0\le N&\le b-a,\label{l:grid:level-bound}\\
 \mathbb E J
 &\le-H(Z\mid G,W)+H(b\mid G,W)
       +q(\log2)\mathbb E(b-a),\label{l:grid:label-charge}\\
 H(b\mid G,W)+H(N\mid G,W)
 &\le2\bigl(1+\mathbb E(b-a)\bigr)\log2.
 \label{l:grid:level-entropies}
\end{align}
\end{lemma}

\begin{proof}
Fix $(g,w)$ with probability $P_{g,w}>0$.  Conditioning gives
$\mu\le\sigma/P_{g,w}$.  Outside one $\mu$-null set, all the cells
$Q_j(s)$ have positive $\mu$-mass: at each level there are only countably
many cells, and a countable union over levels still has zero mass.
We make arbitrary choices on that null set.  On the remaining points,
the score at level $a$ is zero, whereas \eqref{l:grid:cell-mass} bounds
the score at level $j$ by
\[
 \log(1/P_{g,w})-(d-1-q)j\log2-qa\log2.
\]
This upper bound tends to $-\infty$ uniformly in $s$.  Hence a maximum
is attained, and the least maximizing index is measurable because the
scores are measurable and indexed by the integers.  The nonnegative
score at the maximizer also gives
\[
 (d-1-q)b\log2\le\log(1/P_{g,w}).
\]
Averaging proves $\mathbb E b<\infty$.  Since
$\log\mu(Q_b(S))\le0$, we have
$J\le q(b-a)\log2\le d(b-a)$, proving \eqref{l:grid:level-bound}.

For a nonnegative integer variable $V$ of finite mean, relative entropy
to the probability masses $2^{-i-1}$ gives
\[
 H(V)\le(1+\mathbb E V)\log2.
\]
Apply this conditionally to $b-a$ and to $N$.  The value of $a$ is
known from $G$, so $H(b\mid G,W)=H(b-a\mid G,W)$; moreover
$\mathbb E N\le\mathbb E(b-a)$.  This proves
\eqref{l:grid:level-entropies}.  Given $b=j$, the cell label has at
most $(C2^j)^d$ possibilities.  Hence
$H(Q_b(S)\mid G,W,b)\le d\log C+d(\log2)\mathbb E b$, and
$H(G,W,Z)$ is finite.

It remains to account for the selected cell label.  Fix $g,w$ and a
value $b=j$ of positive conditional probability.  Compare the
conditional distribution of $Q_j(S)$ under this selection with the
level-$j$ cell probabilities of $\mu$.  Nonnegative relative entropy
gives
\[
 \mathbb E[\log\mu(Q_j(S))\mid g,w,b=j]
 \le-H(Q_j(S)\mid g,w,b=j).
\]
Averaging this inequality and using the defining formula for $J$ gives
\eqref{l:grid:label-charge}.  This argument conditions on the selected
level; it uses no stopping-time property of that level.
\end{proof}

The next lemma measures the information gained from one block.  Its
posterior may be irregular before the classes indexed by $N$ are formed.
The factor $1/p$ in the state-alphabet cost comes from a $p$th projection
moment.

\Needspace{6\baselineskip}
\begin{lemma}[Information acquired inside a grid cell]
\label{l:grid:block}
Let $S\sim\sigma$.  Let $G$ be countable with $H(G)<\infty$, and suppose
that it determines an integer $a\ge0$ and a level-$a$ cube containing
$S$, with $\mathbb E a<\infty$.  Independently of $(S,G)$, draw a
standard Gaussian $k$-by-$d$ matrix $A$.  Given $(A,AS,G)$, let a
measurable probability kernel produce $W$, with at most $e^m$ possible
values for each $G=g$, where $m\ge0$.  Form $J,N$ from $(G,W)$ as in
Lemma~\ref{l:grid:refinement}.  Then
\begin{equation}\label{l:grid:block-bound}
 I(S;W\mid G)\le Cd+\frac mp+\mathbb E J+H(N\mid G,W).
\end{equation}
\end{lemma}

\begin{proof}
Here $H(G,W)\le H(G)+m<\infty$, so the refinement lemma applies.
We first bound projection densities for each class, using a fresh
Gaussian matrix.  We then compare the actual matrix, selected through
$W$, with that fresh law.

Fix $g,w,n$ of positive probability.  Put
\[
 \alpha=\mathbb P(N=n\mid g,w),\qquad
 \mu_n=\mathcal L(S\mid g,w,n),\qquad \delta=2^{-a}.
\]
A level-$j$ cube of positive $\mu_n$-mass contains a class point with
$J<(n+1)d$.  Since the score at that level is at most $J$, and
$\mu_n\le\mu/\alpha$, it follows that
\begin{equation}\label{l:grid:class-cell}
 \mu_n(Q_j)\le\alpha^{-1}e^{(n+1)d}(2^{a-j})^q
 \qquad(j\ge a).
\end{equation}
For $0<t\le1$, choose $j\ge a$ with
$t\delta/2\le\delta_j\le t\delta$.  A ball of radius $t\delta$
meets at most $C^d$ such cubes: their union lies in a ball of radius
$2t\delta$, and the volume estimate used above cancels the factor
$d^{d/2}$ in the inverse cube volume.  Thus \eqref{l:grid:class-cell}
bounds the mass of this ball by
$C^de^{(n+1)d}\alpha^{-1}t^q$.

Let $V$ be any affine plane of dimension $r\le p-1$.  The projection
of the current level-$a$ cube onto $V$ lies in a plane-ball of radius
$\delta$.  A maximal separated set covers that ball by at most
$(1+2/t)^r$ balls of radius $t\delta$.  Points of the current cube
within distance $t\delta$ of $V$ then lie in the corresponding ambient
balls of radius $2t\delta$.  For $t\le1/2$ apply the ball estimate;
for $t>1/2$ enlarge $C$ and use total mass one.  We obtain
\begin{equation}\label{l:grid:tube}
 \mu_n\{u:\operatorname{dist}(u,V)\le t\delta\}
 \le C^de^{(n+1)d}\alpha^{-1}t^{q-r}
 \qquad(0<t\le1).
\end{equation}
In particular every such plane has zero mass.  The layer-cake formula,
applied to $(\delta/\operatorname{dist}(u,V))^k$, gives
\begin{equation}\label{l:grid:negative-moment}
 \int\operatorname{dist}(u,V)^{-k}\,d\mu_n(u)
 \le\delta^{-k}C^de^{(n+1)d}\alpha^{-1}.
\end{equation}
Indeed the integral after multiplication by $\delta^k$ is at most
$1+Kk/(q-r-k)$ with $K=C^de^{(n+1)d}/\alpha$, and
$q-r-k$ is a positive multiple of $d$.

Fix any $s\in S^{d-1}$ and draw $U_1,\ldots,U_p$ independently from
$\mu_n$.  Let $\Gamma_s$ be the Gram matrix of the vectors $U_i-s$.
Gram--Schmidt expresses its determinant as the product of the squared
distances from $U_i$ to
$s+\operatorname{span}\{U_h-s:h<i\}$.  These planes have dimension at
most $p-1$, and \eqref{l:grid:tube} shows that the distances are positive
almost surely.  Successive integration using
\eqref{l:grid:negative-moment} yields
\begin{equation}\label{l:grid:gram-moment}
 \int(\det\Gamma_s)^{-k/2}\,d\mu_n^{\otimes p}
 \le\bigl(\delta^{-k}C^de^{(n+1)d}\alpha^{-1}\bigr)^p.
\end{equation}
The center $s$ is fixed in this estimate and need not have law $\mu_n$.

We turn this inverse-volume bound into a $p$th moment of the class
projection density.  That moment will control its average logarithm
under the actual matrix law; the relative-entropy comparison will charge
at most $m/p$ for selection through $W$.

Let $\phi_\eta$ be the $N(0,\eta^2I_k)$ density and define
\[
 \rho_\eta(A,y)=\int\phi_\eta(y-Au)\,d\mu_n(u).
\]
Write $\gamma_k$ for standard Gaussian $k$-by-$d$ matrix law.  For a
fresh $\widetilde A\sim\gamma_k$, expansion of the $p$th power gives
\begin{equation}\label{l:grid:mollified-moment}
 \mathbb E_{\widetilde A}
       \rho_\eta(\widetilde A,\widetilde A s)^p
 =(2\pi)^{-kp/2}\int
       \det(\Gamma_s+\eta^2I_p)^{-k/2}\,d\mu_n^{\otimes p}.
\end{equation}
For each row, the vector of projected differences has covariance
$\Gamma_s$; adding the $p$ independent mollifying Gaussians adds
$\eta^2I_p$.  Evaluating its density at zero and multiplying across
independent rows proves the identity.  The right side is bounded by
\eqref{l:grid:gram-moment}.

We next select one density version that is valid on the true-label
diagonal as well as under the actual selected rows.  The projection of
uniform sphere measure onto $k<d$ coordinates has density
\[
 \frac{\Gamma(d/2)}{\pi^{k/2}\Gamma((d-k)/2)}
 (1-\|y\|^2)^{(d-k-2)/2}\mathbf1_{\{\|y\|<1\}}.
\]
For completeness, write a uniform sphere point as a normalized standard
Gaussian.  The squared length of its first $k$ coordinates is the ratio
$X/(X+X')$ for independent chi-square variables of $k$ and $d-k$
degrees of freedom.  Changing variables in their gamma densities gives
the beta density with parameters $k/2,(d-k)/2$; its direction is uniform
and independent.  Polar coordinates give the displayed density.
Rotations and an invertible range map prove absolute continuity for
every rank-$k$ linear map, and the property passes from $\sigma$ to
$\mu_n\ll\sigma$.

Consequently the law of $(\widetilde A,\widetilde A U)$, with
$U\sim\mu_n$ independent of $\widetilde A$, has a density $\rho_0$
relative to $\gamma_k(dA)\,dy$.  The functions $\rho_\eta$ converge to
$\rho_0$ in this joint $L^1$ space.  To see this, for almost every $A$
use continuity of translations in $L^1(dy)$ for its conditional
density and average the Gaussian translations; the error is at most two,
so dominated convergence permits averaging over $A$.  Choose a sequence
$\eta_j\downarrow0$ along which convergence also holds almost everywhere,
and define everywhere
\[
 \rho(A,y)=\liminf_{j\to\infty}\rho_{\eta_j}(A,y).
\]
This is a version of $\rho_0$.  The sequence was chosen without a value
of $s$.  Since \eqref{l:grid:mollified-moment} was bounded for every
fixed $s$, Fatou's lemma now gives, for this same version and for each
fixed $s$,
\begin{equation}\label{l:grid:diagonal-moment}
 \left(\mathbb E_{\widetilde A}
       \rho(\widetilde A,\widetilde A s)^p\right)^{1/p}
 \le\delta^{-k}C^de^{(n+1)d}\alpha^{-1}.
\end{equation}
There are only countably many triples $(g,w,n)$, so this construction
can be made for every occurring triple.  Conditional on such a triple,
the actual matrix law is absolutely continuous with respect to
$\gamma_k$.  Also, for almost every such matrix, the conditional signal
law is absolutely continuous with respect to $\sigma$: conditioning the
pre-block law by the discrete events $W=w,N=n$ only reweights it.
Its image under a full-rank matrix therefore has a Lebesgue density.
Thus the actual joint law of the matrix and its label is absolutely
continuous with respect to $\gamma_k(dA)\,dy$.  In particular the
chosen version $\rho$ is valid in expectations under that actual law.

We have obtained the fresh projection estimate.  To convert it into
information, set $Y=AS$.  Fresh independence and the conditional Markov
property $S\longrightarrow Y\longrightarrow W$ given $(A,G)$ imply
\begin{align}
 I(S;W\mid G)
 &=I(Y;W\mid A,G)-I(A;S\mid G,W)\notag\\
 &=h(Y\mid A,G)-h(Y\mid A,G,W)-I(A;S\mid G,W).
 \label{l:grid:information-identity}
\end{align}
The first line follows by expanding $I(S;A,W\mid G)$ in two orders.
It also gives $I(A;S\mid G,W)\le I(S;W\mid A,G)\le m$.

All the differential entropies in this argument, including after
conditioning on $N$, are finite.  Here are the needed integrability
details.  Because $d-k\ge2$, the coordinate density displayed above is
bounded by a finite constant depending on $d,k$.  For a rank-$k$ matrix
the corresponding bound is multiplied by
$(\det AA^{\mathsf T})^{-1/2}$.  Gram--Schmidt on Gaussian rows writes
this determinant as a product of squared residual lengths, each with a
chi-square law; the absolute logarithm is integrable from its gamma density.
Thus $\mathbb E|\log\det AA^{\mathsf T}|<\infty$.  Given $G=g$, the
signal density is bounded by $1/\mathbb P(G=g)$ relative to $\sigma$,
whose expected logarithm is $H(G)$.  Conditioning further on $(W,N)$
can multiply a conditional label density by at most
$1/\mathbb P(W=w,N=n\mid A,G)$; its expected logarithm is
$H(W,N\mid A,G)\le H(W,N\mid G)<\infty$.  These facts control the
positive parts of the log label densities.  To control their negative
parts, compare with a fixed standard Gaussian label density.  For two
densities $f,g$,
\[
 \int_{\{f<g\}} f\log(g/f)\le 1/e,
\]
because $t\log(1/t)\le1/e$ for $0\le t\le1$.  The remaining Gaussian
logarithm is integrable since $\mathbb E\|Y\|^2=k$.  This proves the
claimed finiteness and justifies the entropy identity.

Let $c_G$ be a specified corner of the current cube, so
$\|S-c_G\|\le\delta=2^{-a}$.  Comparing the conditional label density
with $N(Ac_G,\delta^2I_k)$, and using independence of $A$ to obtain
$\mathbb E[\|A(S-c_G)\|^2\mid G]\le k\delta^2$, gives
\begin{equation}\label{l:grid:upper-entropy}
 h(Y\mid A,G)\le k\mathbb E\log\delta+Ck.
\end{equation}
For the other two terms in \eqref{l:grid:information-identity}, first
condition on $N$.  Differential entropy decreases under this
conditioning, and, since $N$ is determined by $(S,G,W)$,
\[
 I(A;S\mid G,W)=I(A;N\mid G,W)+I(A;S\mid G,W,N).
\]
The latter conditional mutual information is the average divergence
from $\mathcal L(S\mid A,G,W,N)$ to its marginal $\mu_N$.
Data processing under $s\mapsto As$, followed by the definition of
relative entropy for the two label densities, gives
\begin{align*}
 &h(Y\mid A,G,W,N)+I(A;S\mid G,W,N)\\
 &\quad\ge h(Y\mid A,G,W,N)
  +\mathbb E D\!\left(\mathcal L(Y\mid A,G,W,N)
       \,\middle\|\,\rho_{G,W,N}(A,y)\,dy\right)\\
 &\quad=-\mathbb E\log\rho_{G,W,N}(A,AS).
\end{align*}
Together with the preceding conditioning and chain-rule inequalities,
this proves
\begin{equation}\label{l:grid:lower-bracket}
 h(Y\mid A,G,W)+I(A;S\mid G,W)
 \ge-\mathbb E\log\rho_{G,W,N}(A,AS).
\end{equation}
This comparison is between finite quantities: the conditional mutual
information is at most $m$, and the true log density was shown integrable.
The projected log density ratio has integrable negative part by the same
$1/e$ bound and finite positive part by data processing.

If $P\ll P_0$ and $f\ge0$, relative entropy to the probability law
with density proportional to $f^p$ gives
\begin{equation}\label{l:grid:entropy-tilt}
 \mathbb E_P\log f\le
 \frac1p\left(D(P\|P_0)+\log\mathbb E_{P_0}f^p\right).
\end{equation}
Truncating $\log f$ proves this form whenever the left side is defined
and the right side is finite.  In the present application the preceding
paragraph gives integrability.
Apply \eqref{l:grid:entropy-tilt} conditionally on $(s,g,w,n)$ to the
actual matrix law, taking $P_0=\gamma_k$ and
$f(A)=\rho_{g,w,n}(A,As)$.  The value of $n$ is already determined by
$(s,g,w)$, and before $W$ the matrix is independent of $(S,G)$.
The average divergence in this comparison is therefore
\[
 I(A;W\mid S,G)\le H(W\mid G)\le m.
\]
The average of $-\log\alpha$ is $H(N\mid G,W)$.  Using
\eqref{l:grid:diagonal-moment} in \eqref{l:grid:entropy-tilt} and then
averaging gives
\[
 \mathbb E\log\rho_{G,W,N}(A,AS)
 \le-k\mathbb E\log\delta+Cd+\mathbb E[(N+1)d]
       +H(N\mid G,W)+m/p.
\]
Finally $(N+1)d\le J+d$.  Substitute this bound and
\eqref{l:grid:upper-entropy}--\eqref{l:grid:lower-bracket} into
\eqref{l:grid:information-identity} to obtain
\eqref{l:grid:block-bound}.
\end{proof}

\Needspace{6\baselineskip}
\begin{proposition}[Iteration of the grid information balance]
\label{l:grid:iteration}
Let $S\sim\sigma$, and start from a fixed state.  Let $B\ge0$ be an integer.
Consider $B$ blocks, each consisting of $k$ fresh standard Gaussian rows and their
exact labels.  In each block, the outgoing state is a measurable kernel
of the incoming state and that block's data, and its alphabet has
logarithm at most $m\ge0$.  There is a countable analysis history $G_B$ of
finite entropy that contains the final state and satisfies
\begin{equation}\label{l:grid:final-information}
 I(S;G_B)\le C\{d+B(d+m/p)\}.
\end{equation}
\end{proposition}

\begin{proof}
Start with $G_0=Q_0(S)$ and depth $a_0=0$.  Its entropy is at most $Cd$
by the grid count.  After every outgoing state $W$, append the pair
$Z=(b,Q_b(S))$ from Lemma~\ref{l:grid:refinement}.  The enlarged history
contains the outgoing state and a cell containing $S$ at depth $b$.
Its entropy and expected depth remain finite.  Since this refinement
uses only $S$ and completed blocks, the next Gaussian rows remain
independent of $(S,G)$.

For a current history of depth $a$, define
\[
 F(G)=I(S;G)-\bar q(\log2)\mathbb E a.
\]
Because $Z$ is a function of $(S,G,W)$,
\[
 I(S;G,W,Z)-I(S;G)=I(S;W\mid G)+H(Z\mid G,W).
\]
The block estimate, \eqref{l:grid:label-charge}, and
\eqref{l:grid:level-entropies} give
\[
 F(G,W,Z)\le F(G)+Cd+m/p.
\]
Here the term $q\log2$ pays the selected cell mass, and the remaining
$2\log2$ per unit of depth pays both the depth and class entropies.
Thus the entire label $Z$ has been included in the balance.

For each history value $g$, its conditional signal law is supported
on the specified level-$a$ cube.  Relative entropy to $\sigma$ is at
least the negative logarithm of the $\sigma$-mass of that cube: this
follows by writing it as relative entropy to normalized $\sigma$ on
the cube minus the logarithm of the cube mass.  By
\eqref{l:grid:cell-mass},
\[
 I(S;G)\ge(d-1)(\log2)\mathbb E a.
\]
It follows that
$F(G)\ge(1-\bar q/(d-1))I(S;G)$, with a positive coefficient bounded
away from zero for large $d$.  Iterating the upper bound from
$F(G_0)\le Cd$ proves \eqref{l:grid:final-information}.
\end{proof}

\paragraph{Application to precision and bounded stopping.}
\begin{corollary}[Precision under the uniform spherical prior]
\label{l:grid:precision}
Let $M(d)=o(d^2)$ and $0<\epsilon(d)\le1/10$.  For learners in the
finite-state model of Section~\ref{l:model}, uniform-sphere average
angular success at least $2/3$ requires
\[
 T(d)\ge c\,d\log\frac1{\epsilon(d)}
\]
for an absolute $c>0$ and all sufficiently large $d$.  The dimension
threshold may depend on the memory sequence.
\end{corollary}

\begin{proof}
Use the reductions following Lemma~\ref{l:prior-ac} under the uniform
prior.  First fix the shared rule randomness with success greater than
$3/5$, and apply Lemma~\ref{l:borel}.  Then fix the independent
initial state, the finite transition choices, and the outputs assigned
to terminal state/index pairs by averaging.  The resulting deterministic
Borel learner has uniform-prior success at least $3/5$ and the same
state and horizon bounds.

Put $L=\log(1/\epsilon)$, with $0<\epsilon\le1/10$.  The case $T\ge dL$
already has the required order.  In the other case the fixed learner has
at most $(T+1)2^M$ terminal outputs, counting every possible stopping
index.  Angular error at most $\epsilon$ implies Euclidean error at most
$\epsilon$, because $2\sin(\theta/2)\le\theta$.  Applying
\eqref{l:grid:cap} to each output gives
\begin{equation}\label{l:grid:output-count}
 \frac35\le(T+1)2^M\epsilon^{d-1},\qquad
 (d-2)L\le M\log2+O(\log d).
\end{equation}
For the second inequality, take logarithms in the first and use
$T<dL$ and $\log(dL+1)\le\log(d+1)+L$.  Thus $M=o(d^2)$ implies
$L=o(d)$ along these putative short runs, and eventually $T\le d^2$.
This is a consequence of the short-run assumption, not an added
restriction on the accuracy sequence.

Preserve a stopped state together with its terminal index, and use
ignored fresh samples after stopping.  At any later index there are at
most $2^M$ active states and $(T+1)2^M$ tagged terminal states.  Padding
the last block as well therefore uses width at most
\begin{equation}\label{l:grid:padded-width}
 (T+2)2^M,\qquad m=M\log2+\log(T+2)=o(d^2).
\end{equation}
Take $B=\lceil T/k\rceil$ blocks, padding with fewer than $k$ ignored
samples if necessary.  The final augmented history recovers the output.
Under the product of the signal and output marginals, angular success
has probability at most $\epsilon^{d-1}$ by \eqref{l:grid:cap}.
For an event with probabilities $u$ and $v$ under two laws, data
processing to the event gives relative entropy at least
$u\log(1/v)-\log2$: expand binary relative entropy and use binary
entropy at most $\log2$.  Hence
\begin{equation}\label{l:grid:precision-information}
 I(S;G_B)\ge I(S;\widehat S)
 \ge\frac35(d-1)L-\log2.
\end{equation}
Since $m/p=o(d)$, \eqref{l:grid:final-information} and
\eqref{l:grid:precision-information} imply $L\le C(1+B)$.
With $B\le T/k+1$, this gives $T\ge cdL$ whenever $L$ exceeds a fixed
absolute constant.

We finish the bounded range of $L$ with the route's linear sample bound.
Grant the learner all $T<d$ pre-generated rows and exact labels, including
any it did not use.  Their row space $V$ has dimension $T$ almost surely.
The data determine $P_VS$, and conditional on the data the residual
$S-P_VS$ has uniform direction in $V^\perp$, with radius
$\sqrt{1-\|P_VS\|^2}$.  This follows by splitting a normalized standard
Gaussian into its independent row-space and nullspace components.
Its conditional mean is zero and its covariance is
$(1-\|P_VS\|^2)P_{V^\perp}/(d-T)$.  For every unit estimate determined
by the data, also after averaging independent output randomness,
\begin{equation}\label{l:grid:linear-endpoint}
 \mathbb E\langle S,\widehat S\rangle^2
 \le\mathbb E\|P_VS\|^2+\frac1{d-T}
 =\frac Td+\frac1{d-T}.
\end{equation}
On angular success the squared inner product is at least
$\cos^2(1/10)$.  Since $(3/5)\cos^2(1/10)>1/2$, average success
$3/5$ excludes $T\le d/2$ for large $d$.  This linear bound absorbs
the bounded range of $L$ and completes the proof.
\end{proof}
